% Section .............................................................................................................
\section{Aristotle's Potentiality and the Quantum State}
\label{section:aristotle}

% Subsection ..........................................................................................................
\subsection{Heisenberg's Invocation}

In \emph{Physics and Philosophy} (1958), Heisenberg introduced an Aristotelian category to interpret quantum mechanics \cite{heisenberg_1958_physics}:

\begin{quote}
  ``[The probability wave] was a quantitative version of the old concept of \textit{potentia} in Aristotelian philosophy. It introduced something standing in the middle between the idea of an event and the actual event, a strange kind of physical reality just in the middle between possibility and reality.''
\end{quote}

This invocation is often treated as a philosophical metaphor appended to an otherwise self-contained physical theory. We argue it is more: it is a serious ontological proposal that deserves reconstruction and critical evaluation on its own terms.

% Subsection ..........................................................................................................
\subsection{The Aristotelian Framework}

Aristotle's \emph{Metaphysics} \cite{aristotle_1924_metaphysics} distinguishes two fundamental modes of being:

\begin{itemize}
  \item \dunamis\ (potentiality): what a thing can be or do, considered as a real but not yet actualised property. Not mere logical possibility, nor ignorance about an already determinate state, but a genuine ontological category---a real property of things, prior to and distinct from their actual states.
  \item \energeia\ / \entelechia\ (actuality): the realisation of that potentiality; what a thing is in full exercise of its nature.
\end{itemize}

For Aristotle, bronze \emph{is} potentially a statue before any sculptor acts; the seed \emph{is} potentially a tree. 
These are not descriptions of what we happen not to know; they describe what the thing really is, in a mode of being distinct from full actuality.

% -----------------------------------------------------------
% INSERT between the Aristotelian framework paragraph and
% "Application to the Quantum State"
% New subsection on Zubiri
% -----------------------------------------------------------

\subsection{A Note on Zubiri and the Etymology of Being}

The recovery of Aristotelian potentiality for physics is not without
precedent in the philosophical tradition, though it has gone largely
unnoticed in the analytic philosophy of physics.
Xavier Zubiri, in \emph{Sobre la esencia} \cite{zubiri_1962_esencia}
and \emph{Inteligencia sentiente} \cite{zubiri_1980_inteligencia},
developed a sustained critique of the ontological tradition precisely
through attention to the etymology and original sense of the Greek
\textgreek{ὄν} (\emph{on}, the present participle of \emph{einai},
to be).

Zubiri's central contention is that the Latin translation of
\textgreek{ὄν} as \emph{ens}---and the subsequent scholastic
hypostasization of \emph{esse} as an abstract noun---transformed
what was in Aristotle a \emph{dynamic} concept (being as an ongoing
process, a \emph{being-in-act}) into a \emph{static} one (being as
a fixed substantial nature).
The classical ontology of the atom inherits precisely this Latin
distortion: electrons are conceived as \emph{entia} with fully
determinate properties at every instant---a picture of being as
complete, static actuality.

What Zubiri calls \emph{la realidad} is not a fixed substrate of
properties but a \emph{de suyo}---a being-from-itself that is
irreducibly open, whose determinations emerge in relation to other
realities and cannot be read off from an intrinsic nature considered
in isolation.
This is structurally congruent with the potentiality reading: the
electron's \dunamis\ is not a deficient mode of being awaiting
completion, but a genuine and irreducible mode of \emph{de suyo}
that precedes and grounds any particular actualisation.

Zubiri's framework has not been applied to the philosophy of quantum
mechanics, to our knowledge.
We suggest it deserves to be: his critique of the static Latin
\emph{esse} maps precisely onto the failure of the classical atomic
ontology, and his recovery of dynamic \textgreek{ὄν} maps onto the
potentiality reading we are defending.
The convergence is not merely terminological---it reflects a shared
diagnosis that the metaphysical presuppositions inherited from
scholastic substance ontology are inadequate to describe entities
whose mode of being is irreducibly open and relational.

% Subsection ..........................................................................................................
\subsection{Application to the Quantum State}

The mapping is structurally precise. Before a position measurement on an electron in state $\psi$, the electron does not have a definite position we merely fail to know---this is what Bell's theorem and its experimental confirmation rule out \cite{bell_1964_epr, aspect_1982_bell, hensen_2015_loophole}.

\medskip

Rather, the electron has a distribution of \emph{potential} positions, encoded by the so-called wave function $|\psi(x)|^2$, which are real properties of the electron's present state. Measurement does not reveal a pre-existing value; it \emph{actualises} one of the potentialities.

\medskip

On this reading:
\begin{itemize}
  \item $\psi$ is the mathematical representation of the electron's potentialities.
  \item Measurement is the physical process of passage from potentiality to actuality.
  \item The Born rule $|\psi(x)|^2\,\mathrm{d}x$ gives the quantitative structure of the potentiality distribution.
\end{itemize}

This is not the Copenhagen interpretation in its instrumentalist form, which treats $\psi$ as a bookkeeping device for predictions. It is a realist reading: the potentialities are real features of the world, not features of our knowledge.

% Subsection ..........................................................................................................
\subsection{What This Resolves and What It Does Not}

The Aristotelian reading dissolves the puzzle of the classical collapse. The $1s$ electron has no definite trajectory; there is no classical kinematic state to insert into Larmor's formula~(\ref{eq:larmor}). The question ``why does the orbiting electron not radiate?'' does not arise, because its presupposition---that the electron has a definite centripetal acceleration---is false.

What this reading does \emph{not} resolve is the measurement problem in its general form. It relocates the hard question to the actualisation process: what determines which potentiality becomes actual? The major interpretations map partially onto different answers:

\begin{itemize}
  \item \textbf{Copenhagen:} actualisation is a primitive, theory-external fact; quantum mechanics describes potentialities but not the process of actualisation itself.
  \item \textbf{GRW:} spontaneous collapse provides a physical mechanism for actualisation, at the cost of modifying the Schr\"{o}dinger equation \cite{ghirardi_1986_unified}.
  \item \textbf{Everett:} all potentialities are actualised in branching worlds; no collapse occurs \cite{everett_1957_relativestate}.
  \item \textbf{Bohmian mechanics:} particles always have definite positions (full actuality); the pilot wave guides them deterministically \cite{bohm_1952_hiddenvariables}.
\end{itemize}

% Añadir tras el párrafo "We do not adjudicate between these here..."

\phantomsection\label{sec:wigner}

The measurement problem becomes sharper still in Wigner's
thought experiment of 1961 \cite{wigner_1961_friend}.
An observer inside an isolated laboratory---Wigner's
friend---measures the spin of an electron and, from his
perspective, obtains a definite outcome.
Wigner, outside the laboratory, describes the entire system
(friend, apparatus, electron) as a closed quantum system
evolving unitarily, assigning it the entangled state
\begin{equation}
  |\Psi\rangle = \frac{1}{\sqrt{2}}
  \bigl(|\uparrow\rangle_e|\mathrm{saw}\,\uparrow\rangle_F
      + |\downarrow\rangle_e|\mathrm{saw}\,\downarrow\rangle_F\bigr).
  \label{eq:wigner}
\end{equation}
The two observers assign incompatible quantum states to the same
physical system at the same moment.
On the potentiality reading, the paradox sharpens a question that
our framework does not fully answer: is the passage from
potentiality to actuality---from $\delta\acute{\upsilon}\nu\alpha\mu\iota\varsigma$
to $\epsilon\nu\acute{\epsilon}\rho\gamma\epsilon\iota\alpha$---an
objective physical event, or is it indexed to a particular
observer's interaction with the system?
If the former, when exactly does it occur?
If the latter, the potentialities are not fully ontic in the
sense we have been defending.

Frauchiger and Renner \cite{frauchiger_2018_quantum} sharpened
this into a no-go theorem: assuming simultaneously that quantum
mechanics is universal, that different observers' reasoning is
mutually consistent, and that the Born rule applies,
leads to contradiction.
The theorem does not refute any single interpretation but shows
that the three assumptions cannot all be true together---each
major interpretation must abandon at least one.
For the potentiality realist, the theorem suggests that the
ontology of actualisation requires further precision: a
fully satisfactory account must specify the physical conditions
under which joint potentialities become actual, and for whom.

\medskip

We do not adjudicate between these here. Our claim is more modest: the Aristotelian framework provides a more adequate \emph{starting} ontology for quantum mechanics than either classical particle realism or strict instrumentalism, and the collapse calculation is the sharpest way to see why.

\TODO{Reviewer will ask: is potentiality realism distinguishable from an epistemic interpretation? Key argument: potentialities are ontic, not epistemic---they ground the Born-rule probabilities rather than describing ignorance. See \cite{margenau_1954_interpretations} and \cite{bunge_1979_ontology}.}

% -----------------------------------------------------------
% INSERT replaces the \TODO at the end of section 4.4
% -----------------------------------------------------------

A natural objection must be addressed directly: is the potentiality
reading genuinely \emph{ontic}, or is it merely a reformulation of an
epistemic interpretation in Aristotelian vocabulary?
The objection runs as follows---the wavefunction $\psi$ describes our
knowledge of the system; ``potentiality'' is just a philosophically
dressed version of ``probability given our ignorance''; and the
Aristotelian language adds no ontological content.

We resist this reading on two grounds.

First, Bell's theorem \cite{bell_1964_epr} and its experimental
confirmation \cite{aspect_1982_bell,hensen_2015_loophole} rule out
any interpretation on which the wavefunction merely encodes ignorance
about pre-existing local values.
If $|\psi(x)|^2$ described our ignorance of a definite pre-existing
position, the CHSH inequality would not be violated.
It is violated.
The potentialities encoded in $\psi$ cannot therefore be epistemic in
the relevant sense---they are not ignorance about something that is
already actual.

Second, Margenau argued as early as 1954 \cite{margenau_1954_interpretations}
that quantum states describe \emph{latencies}---real dispositional
properties of physical systems that are not reducible to the statistics
of measurement outcomes.
Bunge \cite{bunge_1979_ontology} developed this into a systematic
ontology of propensities, arguing that quantum probabilities are
grounded in real physical tendencies of systems, not in subjective
degrees of belief or objective frequencies over ensembles.
On both accounts, the potentialities are \emph{ontic}: they are
features of the world that ground the Born-rule statistics, not
descriptions of our epistemic state with respect to an already-actual
world.


\footnote{Barandes \cite{barandes_2023_stochastic} shows that the
full empirical content of quantum mechanics can be recovered from
\emph{indivisible stochastic processes}---transition probabilities
satisfying a generalised Markov condition---without positing a
wavefunction or Hilbert space as primitive ontology. Potentialities
become transition probabilities between configurations, arguably
closer to the Aristotelian notion of \dunamis\ as process than the
static wavefunction picture. The Barandes framework is too recent to evaluate
fully, but illustrates that the ontological layer of quantum
mechanics remains genuinely open}

The potentiality reading we are defending is therefore not epistemic
Copenhagen in disguise.
It is a realist position that takes the wavefunction to represent
real features of the world---features that are genuinely potential
rather than actual, in the Aristotelian sense, prior to measurement.
The difference from standard collapse interpretations is that we do
not treat potentiality as a temporary state of ignorance to be
resolved by measurement into something fully actual and classical.
Potentiality is a permanent and irreducible mode of being for
quantum systems between interactions---it is what they \emph{are},
not what we \emph{don't know} about them.


% Como \footnote{} al final del párrafo que lista Copenhagen/GRW/Everett/Bohm